% ================================================================
% COMMANDS
% Perintah (command) khusus yang dapat dipakai ulang.
% Semua nilai dokumen berasal dari preamble/metadata.tex (\doc*).
% ================================================================

\newcommand{\HRule}{\rule{\linewidth}{0.5mm}}

% ================================================================
%% FIGURE
%% Sisipkan gambar dengan caption dan label dalam satu perintah.
%%
%% Pemakaian:
%%   \Figure{figures/logo_institusi.png}{Keterangan gambar.}{fig:logo}
%%
%% Argumen: [1] file gambar, [2] caption, [3] label.
%% ================================================================

\NewDocumentCommand{\Figure}{m m m}{
  \begin{figure}[ht]
  \centering
  \includegraphics[width=0.8\linewidth]{#1}
  \caption{#2}
  \label{#3}
  \end{figure}
}

% ================================================================
%% TABLE
%% Sisipkan tabel dengan caption dan label dalam satu perintah.
%% Isi tabel ditulis dalam argumen kedua (gaya booktabs).
%%
%% Pemakaian:
%%   \Table{tab:contoh}{%
%%     \begin{tabular}{ll}
%%       \toprule
%%       A & B \\
%%       \bottomrule
%%     \end{tabular}
%%   }{Keterangan tabel.}
%%
%% Argumen: [1] label, [2] isi tabel, [3] caption.
%% ================================================================

\NewDocumentCommand{\Table}{m m m}{
  \begin{table}[ht]
  \centering
  #2
  \caption{#3}
  \label{#1}
  \end{table}
}

% ================================================================
%% EQUATION
%% Sisipkan persamaan bernomor dengan label.
%%
%% Pemakaian:
%%   \Equation{eq:contoh}{E = mc^2}
%%
%% Argumen: [1] label, [2] isi persamaan.
%% ================================================================

\NewDocumentCommand{\Equation}{m m}{
  \begin{equation}
  #2
  \label{#1}
  \end{equation}
}

% ================================================================
%% APPENDIX
%% Mulai bagian lampiran. Membaca label optional untuk penamaan.
%%
%% Pemakaian:
%%   \Appendix{Lampiran}        % bagian lampiran, judul "Lampiran"
%%   \Appendix*                 % lampiran tanpa nomor bab baru
%% ================================================================

\NewDocumentCommand{\Appendix}{s m}{
  \IfBooleanTF{#1}
    {\chapter*{#2}}
    {\chapter{#2}}
}

% ================================================================
%% SOURCE
%% Catatan sumber untuk gambar/tabel (ditempatkan setelah lingkungan).
%%
%% Pemakaian:
%%   \Figure{figures/x.png}{Keterangan.}{fig:x}
%%   \Source{Diadaptasi dari: Author (2020).}
%% ================================================================

\NewDocumentCommand{\Source}{m}{
  \par\smallskip
  \begin{flushleft}
  \small #1
  \end{flushleft}
}

% ================================================================
%% MAKECOVER
%% Halaman sampul (titlepage) yang dapat dipakai ulang.
%% Membaca seluruh identitas dari metadata.tex.
%%
%% Pemakaian:
%%   \makecover                          % judul dari \docTitle
%%   \makecover[Judul Pengganti]         % override judul
%%   \makecover[][Skripsi]               % jenis dokumen (proposal/skripsi/tesis/disertasi/laporan)
%%
%% Contoh jenis dokumen: Proposal, Skripsi, Tesis, Disertasi,
%% Laporan Akhir, dsb. Biarkan argumen ke-2 kosong untuk mengabaikan.
%% ================================================================

\makeatletter

\NewDocumentCommand{\makecover}{ o o }{%
  \begin{titlepage}

  \thispagestyle{empty}

  \begin{center}

  % JUDUL
  {\bfseries
  \fontsize{16pt}{20pt}\selectfont
  \MakeUppercase{\IfValueTF{#1}{#1}{\docTitle}}
  }

  % OPTIONAL: SUBJUDUL
  \ifx\docSubtitle\@empty\relax\else
  {\fontsize{14pt}{18pt}\selectfont
  \docSubtitle
  }
  \fi

  % OPTIONAL: JENIS DOKUMEN
  \IfValueT{#2}{%
  \vspace{1cm}
  {\fontsize{12pt}{18pt}\selectfont
  #2
  }
  }

  \vspace{1.5cm}

  % LOGO (path dari \docLogo)
  \includegraphics[width=4cm,height=4cm,keepaspectratio]{\docLogo}

  \vspace{2cm}

  % PENULIS
  {\fontsize{12pt}{18pt}\selectfont
  Disusun oleh
  }

  \vspace{0.5cm}

  {\bfseries
  \fontsize{12pt}{18pt}\selectfont
  \docAuthor
  }

  \vspace{1.5cm}

  % PEMBIMBING
  {\fontsize{12pt}{18pt}\selectfont
  Telah diperiksa dan disetujui oleh : \\
  }

  \vspace{0.5cm}

  {\bfseries
  \fontsize{12pt}{18pt}\selectfont
  \docSupervisor
  }

  \vfill

  % INSTITUSI
  {\bfseries
  \fontsize{12pt}{18pt}\selectfont

  \docProgram

  %\docDepartment

  %\docFaculty

  \docHospital

  \docLocation

  \docYear

  }

  \end{center}

  \end{titlepage}

  \clearpage
}

\makeatother
